% METHOD SUPPLEMENT START
\subsection{Navigation Objective and Belief Parameterization}
\label{app:belief-parameterization}

At the system level, task-relevant state can include object placement, spatial
configuration, and accessibility. The portable-object task isolates hidden world
evolution from these broader sources of change. An environment contains
navigable space $\mathcal X$, persistent entities $\mathcal O$, and candidate
states $\mathcal C_o$. At decision step $j$, the agent selects a candidate
viewpoint or a low-level navigation action, receives observation $z_j$, and may
terminate with \texttt{STOP} or \texttt{NOT\_FOUND}. Its objective balances
successful navigation, path length $L$, and inspection count:
\begin{equation}
\pi^*=\arg\max_{\pi}\;\mathbb E_{\pi}\!\left[
\mathbb I(\mathrm{success})-\lambda_d L-\lambda_n N_{\mathrm{inspect}}
\right].
\label{eq:objective}
\end{equation}

\paragraph{Causal memory interface.}
The persistent semantic scene is materialized from versioned deltas:
\begin{equation}
\mathcal M_{\leq t}=\operatorname{Materialize}(\Delta\mathcal M_0,\ldots,\Delta\mathcal M_t)
=(\mathcal G_t,\mathcal V_t,\mathcal B_t),
\label{eq:memory}
\end{equation}
where $\mathcal G_t$ stores semantic--spatial relations, $\mathcal V_t$ stores
time-valid entity versions, and $\mathcal B_t$ stores observation provenance.
The target query returns the most recent $K$ causal events:
\begin{equation}
H_o=\{e_k=(o_k,s_k,t_k,z_k,c_k,v_k)\}_{k=1}^{K}.
\label{eq:history}
\end{equation}
Context events may share an entity, receptacle, region, or time window with the
target. These relations provide context without revealing hidden activities or
transition labels.

\paragraph{Continuous-time event encoding.}
Each event token combines entity, state, evidence, confidence, and visibility
with elapsed and calendar time:
\begin{align}
\mathbf x_k={}&E_o(o_k)+E_s(s_k)+E_z(z_k)+E_c(c_k)+E_v(v_k)\nonumber\\
&+\phi_{\Delta}(t_q-t_k)+\phi_{\mathrm{cal}}(t_k),
\label{eq:eventtoken}
\end{align}
where $\phi_{\Delta}$ includes $\log(1+t_q-t_k)$ and
$\phi_{\mathrm{cal}}$ contains sinusoidal hour-of-day and day-of-week features.
A learned query token specifies the target, query time, and last positive state:
\begin{equation}
(\mathbf h_1,\ldots,\mathbf h_K,\mathbf h_q)
=\operatorname{CTTransformer}(\mathbf x_1,\ldots,\mathbf x_K,\mathbf x_q).
\end{equation}
This representation models irregular gaps, repeated negative evidence, and event
order without observing the activity responsible for a hidden change.

\paragraph{Persistence and relocation heads.}
For last positive state $s_{\mathrm{last}}$, the persistence branch is
\begin{equation}
\rho_q=P(s_q^o=s_{\mathrm{last}}\mid H_o,t_q)
=\sigma(\mathbf w_{\rho}^{\top}\mathbf h_q).
\label{eq:persistence}
\end{equation}
For candidate representation $\mathbf g_i$, a shared pointer head computes
\begin{equation}
a_i=\frac{(W_q\mathbf h_q)^\top(W_c\mathbf g_i)}{\sqrt d}
+\psi(\mathbf h_q,\mathbf g_i),\qquad
\pi_q(i)=\frac{\exp a_i}{\sum_{j\neq s_{\mathrm{last}}}\exp a_j},
\label{eq:pointer}
\end{equation}
for $i\neq s_{\mathrm{last}}$. This defines the relocation distribution in
\cref{eq:prior}; end-to-end state likelihood training makes the factorization an
inductive bias rather than an oracle change label.

\paragraph{Negative observation likelihood.}
For an inspection at $r$, negative evidence is visibility-qualified:
\begin{equation}
P(z=\mathrm{notseen}\mid s_q^o=i,c,v)=
\begin{cases}
1-cv, & i=r,\\
1, & i\neq r.
\end{cases}
\label{eq:negative}
\end{equation}
The inspection, visibility assessment, and posterior are stored as a new causal
memory event. Only sufficiently tested hypotheses are removed; the belief update
and spatial interfaces remain typed and verifiable even when the frozen
high-level controller expresses goals and tool calls in language.
